% GENERATED FILE. Edit canonical lessons, then run npm run generate:books.
% canonical-source-hash: fnv1a64:078fb4a1154da5a8
% canonical-lessons: KA-C511-tagulu, KA-C511-mamparu, KA-C511-ulbanisu, KA-C511-rogalaksana, KA-C511-lasike, KA-R511-review-505-511

\chapter{Catching Something}
\label{ch:ka-catching-something}
\hlchaptermodality{511}

\begin{chapteropening}
\textbf{By the end of this chapter:} \emph{I can say to catch, drowsiness, to get worse, a symptom and a vaccine.}

Complete the last lesson of chapter 511: I can say to catch, drowsiness, to get worse, a symptom and a vaccine.
\end{chapteropening}

\section[\texorpdfstring{\kn{ತಗುಲು} (tagulu)}{ತಗುಲು (tagulu)}]{\kn{ತಗುಲು} (tagulu) — to catch (an illness); to touch}

\label{lesson:KA-C511-tagulu}



\begin{quote}
Before the new one: say the Kannada for the gums, then the Kannada for the jaw.
\end{quote}

\subsection*{You'll want to know: \kn{ತಗುಲು}}
\textbf{\kn{ತಗುಲು}} — \emph{tagulu} — \textquotedblleft{}to catch (an illness); to touch\textquotedblright{}.

\begin{quote}
\kn{ಮಳೆಯಲ್ಲಿ} \kn{ನೆನೆದಿದ್ದರಿಂದ} \kn{ಮಗುವಿಗೆ} \kn{ಶೀತ} \kn{ತಗುಲಿತು}.
\end{quote}

\begin{quote}
\emph{maḷeyalli nenediddarinda maguvige śīta tagulitu.}
\end{quote}

\begin{quote}
Because the child got soaked in the rain, it caught a cold.
\end{quote}

Past (it …): \kn{ತಗುಲಿತು}, \emph{tagulitu}.

The one who falls ill takes -\kn{ಗೆ}: \kn{ಅವನಿಗೆ} \kn{ಸೋಂಕು} \kn{ತಗುಲಿತು}. Costs too: \kn{ಖರ್ಚು} \kn{ತಗುಲುತ್ತದೆ}, it will cost.

\begin{cousinweb}
\textbf{\kn{ತಗುಲು}} (\emph{tagulu}) is native Dravidian: DEDR 3004 lists Kannada \emph{tagalu}, \emph{tagilu} and \emph{tagulu}, \textquotedblleft{}to come in contact with, touch, hit\textquotedblright{}.
\end{cousinweb}

\begin{grammarlens}[title={What you've built}]
One more word for describing what you lived and felt.
\end{grammarlens}

\subsection*{Your turn}
\begin{itemize}
\raggedright
  \item \emph{Say it:} \emph{tagulu}
  \item \emph{Say it:} \emph{tagulu}, once more
  \item \emph{Say it:} \emph{maḷeyalli nenediddarinda maguvige śīta tagulitu.} — Because the child got soaked in the rain, it caught a cold.
  \item \emph{Recall:} say the Kannada for the gums, then the Kannada for the jaw, then say \emph{tagulu} again
\end{itemize}

\begin{tcolorbox}[breakable,colback=teal!4,colframe=teal!35!black,title={Before you move on}]
What does \kn{ತಗುಲು} mean? (\textquotedblleft{}To catch (an illness); to touch\textquotedblright{}.) Say it once more.
\end{tcolorbox}
\section[\texorpdfstring{\kn{ಮಂಪರು} (mamparu)}{ಮಂಪರು (mamparu)}]{\kn{ಮಂಪರು} (mamparu) — drowsiness; grogginess}

\label{lesson:KA-C511-mamparu}



\begin{quote}
Before the new one: say the Kannada for the jaw, then the Kannada for to catch.
\end{quote}

\subsection*{You'll want to know: \kn{ಮಂಪರು}}
\textbf{\kn{ಮಂಪರು}} — \emph{mamparu} — \textquotedblleft{}drowsiness; grogginess\textquotedblright{}.

\begin{quote}
\kn{ಮಾತ್ರೆ} \kn{ತೆಗೆದುಕೊಂಡ} \kn{ಮೇಲೆ} \kn{ನನಗೆ} \kn{ಮಂಪರು} \kn{ಬಂದಿತು}.
\end{quote}

\begin{quote}
\emph{mātre tegedukoṇḍa mēle nanage mamparu banditu.}
\end{quote}

\begin{quote}
After taking the tablet I felt groggy.
\end{quote}

The haze of half-sleep, after medicine or a long night.

\begin{cousinweb}
\textbf{\kn{ಮಂಪರು}} (\emph{mamparu}) is an everyday Kannada word for a drowsy haze; its origin is not settled.
\end{cousinweb}

\begin{grammarlens}[title={What you've built}]
One more word for describing what you lived and felt.
\end{grammarlens}

\subsection*{Your turn}
\begin{itemize}
\raggedright
  \item \emph{Say it:} \emph{mamparu}
  \item \emph{Say it:} \emph{mamparu}, once more
  \item \emph{Say it:} \emph{mātre tegedukoṇḍa mēle nanage mamparu banditu.} — After taking the tablet I felt groggy.
  \item \emph{Recall:} say the Kannada for the jaw, then the Kannada for to catch, then say \emph{mamparu} again
\end{itemize}

\begin{tcolorbox}[breakable,colback=teal!4,colframe=teal!35!black,title={Before you move on}]
What does \kn{ಮಂಪರು} mean? (\textquotedblleft{}Drowsiness; grogginess\textquotedblright{}.) Say it once more.
\end{tcolorbox}
\section[\texorpdfstring{\kn{ಉಲ್ಬಣಿಸು} (ulbaṇisu)}{ಉಲ್ಬಣಿಸು (ulbaṇisu)}]{\kn{ಉಲ್ಬಣಿಸು} (ulbaṇisu) — to get worse; to flare up}

\label{lesson:KA-C511-ulbanisu}



\begin{quote}
Before the new one: say the Kannada for to catch, then the Kannada for drowsiness.
\end{quote}

\subsection*{You'll want to know: \kn{ಉಲ್ಬಣಿಸು}}
\textbf{\kn{ಉಲ್ಬಣಿಸು}} — \emph{ulbaṇisu} — \textquotedblleft{}to get worse; to flare up\textquotedblright{}.

\begin{quote}
\kn{ಚಳಿಯಲ್ಲಿ} \kn{ಅಜ್ಜನ} \kn{ಕೆಮ್ಮು} \kn{ಉಲ್ಬಣಿಸಿತು}.
\end{quote}

\begin{quote}
\emph{caḷiyalli ajjana kemmu ulbaṇisitu.}
\end{quote}

\begin{quote}
In the cold, grandfather's cough flared up.
\end{quote}

Past (it …): \kn{ಉಲ್ಬಣಿಸಿತು}, \emph{ulbaṇisitu}.

Of an illness, a quarrel or a problem that grows worse.

\begin{cousinweb}
\textbf{\kn{ಉಲ್ಬಣಿಸು}} (\emph{ulbaṇisu}) is built on Sanskrit \emph{ulbaṇa}, \textquotedblleft{}superfluous, abundant, excessive, strong\textquotedblright{} (Monier-Williams), with the verb ending -\emph{isu}: to become too much.
\end{cousinweb}

\begin{grammarlens}[title={What you've built}]
One more word for describing what you lived and felt.
\end{grammarlens}

\subsection*{Your turn}
\begin{itemize}
\raggedright
  \item \emph{Say it:} \emph{ulbaṇisu}
  \item \emph{Say it:} \emph{ulbaṇisu}, once more
  \item \emph{Say it:} \emph{caḷiyalli ajjana kemmu ulbaṇisitu.} — In the cold, grandfather's cough flared up.
  \item \emph{Recall:} say the Kannada for to catch, then the Kannada for drowsiness, then say \emph{ulbaṇisu} again
\end{itemize}

\begin{tcolorbox}[breakable,colback=teal!4,colframe=teal!35!black,title={Before you move on}]
What does \kn{ಉಲ್ಬಣಿಸು} mean? (\textquotedblleft{}To get worse; to flare up\textquotedblright{}.) Say it once more.
\end{tcolorbox}
\section[\texorpdfstring{\kn{ರೋಗಲಕ್ಷಣ} (rōgalakṣaṇa)}{ರೋಗಲಕ್ಷಣ (rōgalakṣaṇa)}]{\kn{ರೋಗಲಕ್ಷಣ} (rōgalakṣaṇa) — a symptom}

\label{lesson:KA-C511-rogalaksana}



\begin{quote}
Before the new one: say the Kannada for drowsiness, then the Kannada for to get worse.
\end{quote}

\subsection*{You'll want to know: \kn{ರೋಗಲಕ್ಷಣ}}
\textbf{\kn{ರೋಗಲಕ್ಷಣ}} — \emph{rōgalakṣaṇa} — \textquotedblleft{}a symptom\textquotedblright{}.

\begin{quote}
\kn{ಜ್ವರ} \kn{ಮತ್ತು} \kn{ಕೆಮ್ಮು} \kn{ಈ} \kn{ಕಾಯಿಲೆಯ} \kn{ರೋಗಲಕ್ಷಣಗಳು}.
\end{quote}

\begin{quote}
\emph{jvara mattu kemmu ī kāyileya rōgalakṣaṇagaḷu.}
\end{quote}

\begin{quote}
Fever and a cough are the symptoms of this illness.
\end{quote}

Plural: \kn{ರೋಗಲಕ್ಷಣಗಳು}, \emph{rōgalakṣaṇagaḷu}.

\kn{ರೋಗ}, disease + \kn{ಲಕ್ಷಣ}, a mark.

\begin{cousinweb}
\textbf{\kn{ರೋಗಲಕ್ಷಣ}} (\emph{rōgalakṣaṇa}) is a learned compound built from Sanskrit \emph{rōga}, \textquotedblleft{}disease\textquotedblright{} (CDIAL 10829), and \emph{lakṣaṇa}, \textquotedblleft{}a line, a mark\textquotedblright{} (CDIAL 10882).
\end{cousinweb}

\begin{grammarlens}[title={What you've built}]
One more word for describing what you lived and felt.
\end{grammarlens}

\subsection*{Your turn}
\begin{itemize}
\raggedright
  \item \emph{Say it:} \emph{rōgalakṣaṇa}
  \item \emph{Say it:} \emph{rōgalakṣaṇa}, once more
  \item \emph{Say it:} \emph{jvara mattu kemmu ī kāyileya rōgalakṣaṇagaḷu.} — Fever and a cough are the symptoms of this illness.
  \item \emph{Recall:} say the Kannada for drowsiness, then the Kannada for to get worse, then say \emph{rōgalakṣaṇa} again
\end{itemize}

\begin{tcolorbox}[breakable,colback=teal!4,colframe=teal!35!black,title={Before you move on}]
What does \kn{ರೋಗಲಕ್ಷಣ} mean? (\textquotedblleft{}A symptom\textquotedblright{}.) Say it once more.
\end{tcolorbox}
\section[\texorpdfstring{\kn{ಲಸಿಕೆ} (lasike)}{ಲಸಿಕೆ (lasike)}]{\kn{ಲಸಿಕೆ} (lasike) — a vaccine}

\label{lesson:KA-C511-lasike}



\begin{quote}
Before the new one: say the Kannada for to get worse, then the Kannada for a symptom.
\end{quote}

\subsection*{You'll want to know: \kn{ಲಸಿಕೆ}}
\textbf{\kn{ಲಸಿಕೆ}} — \emph{lasike} — \textquotedblleft{}a vaccine\textquotedblright{}.

\begin{quote}
\kn{ಮಕ್ಕಳಿಗೆ} \kn{ಶಾಲೆಯಲ್ಲಿ} \kn{ಲಸಿಕೆ} \kn{ಹಾಕಿದರು}.
\end{quote}

\begin{quote}
\emph{makkaḷige śāleyalli lasike hākidaru.}
\end{quote}

\begin{quote}
The children were vaccinated at school.
\end{quote}

\kn{ಲಸಿಕೆ} \kn{ಹಾಕು}, to vaccinate.

\begin{cousinweb}
\textbf{\kn{ಲಸಿಕೆ}} (\emph{lasike}), the Kannada word for a vaccine, is usually traced to Sanskrit \emph{lasīkā}, which Monier-Williams glosses as \textquotedblleft{}watery humour in the body, lymph, serum\textquotedblright{} (Monier-Williams); how it came to name a vaccine is not settled.
\end{cousinweb}

\begin{grammarlens}[title={What you've built}]
One more word for describing what you lived and felt.
\end{grammarlens}

\subsection*{Your turn}
\begin{itemize}
\raggedright
  \item \emph{Say it:} \emph{lasike}
  \item \emph{Say it:} \emph{lasike}, once more
  \item \emph{Say it:} \emph{makkaḷige śāleyalli lasike hākidaru.} — The children were vaccinated at school.
  \item \emph{Recall:} say the Kannada for to get worse, then the Kannada for a symptom, then say \emph{lasike} again
\end{itemize}

\begin{tcolorbox}[breakable,colback=teal!4,colframe=teal!35!black,title={Before you move on}]
What does \kn{ಲಸಿಕೆ} mean? (\textquotedblleft{}A vaccine\textquotedblright{}.) Say it once more.
\end{tcolorbox}
\section[(dialogue)]{Looking back: from Effort and Neglect to Catching Something}

\label{lesson:KA-R511-review-505-511}



\begin{quote}
Say the Kannada for dedication and neglect.
\end{quote}

\subsection*{Your turn}
Say these aloud:

\begin{itemize}
\raggedright
  \item dedication, neglect, an addiction, to strive, to beg
  \item to cover oneself, to talk in sleep, to roll over, to stretch, an eyelid
  \item to belch, to blink, to pant, hiccups, asthma
  \item the brain, a lung, a kidney, the intestine, a muscle
  \item a nerve, a rib, the spine, an arm, to limp
  \item to ache, numbness, to swell, the gums, the jaw
  \item to catch, drowsiness, to get worse, a symptom, a vaccine
\end{itemize}

\begin{tcolorbox}[breakable,colback=teal!4,colframe=teal!35!black,title={Before you move on}]
Dedication? (\textbf{\kn{ಶ್ರದ್ಧೆ}}, \emph{śraddhe}.) Neglect? (\textbf{\kn{ನಿರ್ಲಕ್ಷ್ಯ}}, \emph{nirlakṣya}.) An addiction? (\textbf{\kn{ಚಟ}}, \emph{caṭa}.) To strive? (\textbf{\kn{ಶ್ರಮಿಸು}}, \emph{śramisu}.) To beg? (\textbf{\kn{ಬೇಡು}}, \emph{bēḍu}.) To cover oneself? (\textbf{\kn{ಹೊದೆ}}, \emph{hode}.) To talk in sleep? (\textbf{\kn{ಕನವರಿಸು}}, \emph{kanavarisu}.) To roll over? (\textbf{\kn{ಹೊರಳು}}, \emph{horaḷu}.) To stretch? (\textbf{\kn{ಮೈಮುರಿ}}, \emph{maimuri}.) An eyelid? (\textbf{\kn{ರೆಪ್ಪೆ}}, \emph{reppe}.) To belch? (\textbf{\kn{ತೇಗು}}, \emph{tēgu}.) To blink? (\textbf{\kn{ಮಿಟುಕಿಸು}}, \emph{miṭukisu}.) To pant? (\textbf{\kn{ಏದುಸಿರು} \kn{ಬಿಡು}}, \emph{ēdusiru biḍu}.) Hiccups? (\textbf{\kn{ಬಿಕ್ಕಳಿಕೆ}}, \emph{bikkaḷike}.) Asthma? (\textbf{\kn{ಉಬ್ಬಸ}}, \emph{ubbasa}.) The brain? (\textbf{\kn{ಮೆದುಳು}}, \emph{meduḷu}.) A lung? (\textbf{\kn{ಶ್ವಾಸಕೋಶ}}, \emph{śvāsakōśa}.) A kidney? (\textbf{\kn{ಮೂತ್ರಪಿಂಡ}}, \emph{mūtrapiṇḍa}.) The intestine? (\textbf{\kn{ಕರುಳು}}, \emph{karuḷu}.) A muscle? (\textbf{\kn{ಸ್ನಾಯು}}, \emph{snāyu}.) A nerve? (\textbf{\kn{ನರ}}, \emph{nara}.) A rib? (\textbf{\kn{ಪಕ್ಕೆಲುಬು}}, \emph{pakkelubu}.) The spine? (\textbf{\kn{ಬೆನ್ನುಮೂಳೆ}}, \emph{bennumūḷe}.) An arm? (\textbf{\kn{ತೋಳು}}, \emph{tōḷu}.) To limp? (\textbf{\kn{ಕುಂಟು}}, \emph{kuṇṭu}.) To ache? (\textbf{\kn{ನೋಯು}}, \emph{nōyu}.) Numbness? (\textbf{\kn{ಜೋಮು}}, \emph{jōmu}.) To swell? (\textbf{\kn{ಉಬ್ಬು}}, \emph{ubbu}.) The gums? (\textbf{\kn{ಒಸಡು}}, \emph{osaḍu}.) The jaw? (\textbf{\kn{ದವಡೆ}}, \emph{davaḍe}.) To catch? (\textbf{\kn{ತಗುಲು}}, \emph{tagulu}.) Drowsiness? (\textbf{\kn{ಮಂಪರು}}, \emph{mamparu}.) To get worse? (\textbf{\kn{ಉಲ್ಬಣಿಸು}}, \emph{ulbaṇisu}.) A symptom? (\textbf{\kn{ರೋಗಲಕ್ಷಣ}}, \emph{rōgalakṣaṇa}.) A vaccine? (\textbf{\kn{ಲಸಿಕೆ}}, \emph{lasike}.)
\end{tcolorbox}
